% CCSS 7.G.A.3 — Cross-Sections of Three-Dimensional Figures
\section{Cross-Sections of 3D Figures}

\begin{learningGoals}
\begin{itemize}
    \item Describe the 2D shape formed by slicing a 3D figure with a plane
    \item Identify cross-sections of prisms and pyramids
\end{itemize}
\end{learningGoals}

\begin{conceptBox}{What Is a Cross-Section?}
A \textbf{cross-section} is the flat shape you see when you slice through a 3D figure with a plane — like cutting a loaf of bread.
\begin{itemize}[leftmargin=*]
    \item \textbf{Right rectangular prism} (box): slices can be rectangles, squares, or triangles, depending on the angle of the cut.
    \item \textbf{Right rectangular pyramid}: slices can be triangles, trapezoids, or rectangles.
\end{itemize}
The shape depends on \textbf{where} and \textbf{at what angle} you cut.
\end{conceptBox}

\begin{workedExample}{Slicing a Rectangular Prism}
Imagine slicing a box of cereal:
\begin{itemize}[leftmargin=*]
    \item \textbf{Horizontal cut} (parallel to the base): rectangle
    \item \textbf{Vertical cut} (parallel to one face): rectangle
    \item \textbf{Diagonal cut} (corner to corner): triangle or other polygon
\end{itemize}
\answerHighlight{The cross-section shape depends on the angle of the slice.}
\end{workedExample}

\mascotSays{Every horizontal slice of a prism looks exactly like the base. Cut parallel to the base and you always get the same shape!}

\begin{practiceBox}{Cross-Sections of 3D Figures}
\resetProblems

\prob You slice a rectangular prism with a cut parallel to its base. What shape is the cross-section?
\answerExplain{Rectangle}{A horizontal cut parallel to the base of a rectangular prism always produces a rectangle the same size and shape as the base.}

\prob You slice a rectangular pyramid with a horizontal cut parallel to its base, halfway up. What shape do you get?
\answerExplain{A smaller rectangle}{A horizontal cut of a rectangular pyramid parallel to the base gives a rectangle. It is smaller than the base because the pyramid narrows toward the top.}

\prob You cut a cube with a diagonal slice from one edge to the opposite edge. What shape could the cross-section be?
\answerExplain{A rectangle or a triangle}{Depending on the exact angle, a diagonal slice through a cube can create a rectangle, a triangle, or even a hexagon. One common diagonal gives a rectangle.}

\trueOrFalse{Every cross-section of a sphere is a circle.}
\answerTF{True}

\prob A cone is sliced with a horizontal cut parallel to its base. What shape is the cross-section?
\answerExplain{A circle}{A cone has a circular base. A horizontal cut parallel to the base is a smaller circle. The higher you cut, the smaller the circle.}

\end{practiceBox}
